\chapter{Introducción}

Argentina atraviesa una crisis profunda de credibilidad informativa. Según el Digital News Report 2024 del Reuters Institute \parencite{Newman2024}, solo el 30\% de la población confía en los medios de comunicación, el nivel más bajo de América Latina, mientras que el interés en noticias cayó del 77\% en 2017 al 45\% en 2024. Este escenario se desarrolla sobre un entorno de consumo masivo: 31,3 millones de argentinos son usuarios activos de redes sociales, con WhatsApp como principal canal de difusión informativa, utilizado por el 93\% de los usuarios de internet. En este contexto, la desinformación circula a escala y velocidad que hacen imposible su contención manual.

Las soluciones existentes presentan limitaciones estructurales. Chequeado.com, referente local de verificación, opera de forma manual y solo puede cubrir un subconjunto acotado de afirmaciones. Las herramientas comerciales como Cyabra están orientadas a gobiernos y redacciones profesionales, con costos inaccesibles para el ciudadano común. No existe hoy una herramienta gratuita, automática y orientada al contexto argentino que asista a la ciudadanía en tiempo real.

El presente trabajo propone un sistema de detección automática de desinformación en redes sociales y medios digitales, orientado a ciudadanos argentinos mayores de 16 años y a periodistas y editores que necesitan evaluar la confiabilidad de fuentes. La solución se compone de cuatro módulos integrados en una extensión de Google Chrome con panel web.

\section{Objetivos}

\subsection{Objetivo general}

Desarrollar un servicio de inteligencia artificial para facilitar la identificación de contenido desinformativo en redes sociales y medios digitales argentinos.

\subsection{Objetivos específicos}

\begin{enumerate}
    \item \textbf{Modelo de IA}: Entrenar un modelo de inteligencia artificial para clasificar desinformación en español.
    \item \textbf{Credibilidad de fuente}: Evaluar la credibilidad de la fuente mediante análisis de metadatos.
    \item \textbf{Contraste semántico}: Contrastar el contenido contra fuentes confiables usando similitud semántica.
    \item \textbf{Extensión web + dashboard}: Integrar los módulos en una extensión de Chrome con dashboard web.
\end{enumerate}

\section{Alcance}

El alcance comprende el desarrollo de un prototipo funcional de una aplicación web: una extensión de Google Chrome y un panel web. La solución está orientada a ciudadanos argentinos que consumen noticias en redes sociales y medios digitales, y a periodistas y editores como herramienta de validación. La interfaz está en español con uso exclusivo para Argentina. Los usuarios acceden instalando la extensión en Chrome desde computadora de escritorio. El sistema analiza textos en medios digitales con foco en política y economía, y devuelve un score de probabilidad acompañado de evidencia (links a fuentes que corroboran o contradicen la información).

\subsection{Características del prototipo funcional}

\begin{itemize}
    \item Extensión de Google Chrome instalable desde el repositorio o desde la Chrome Web Store.
    \item Panel web complementario para visualización de análisis históricos y reportes.
    \item Análisis de contenido textual en Twitter/X, Instagram, Facebook, Infobae y Clarín.
    \item Clasificación automática mediante modelo de NLP entrenado en español (BETO o XLM-RoBERTa).
    \item Evaluación de credibilidad de fuente basada en metadatos.
    \item Contraste semántico del contenido contra fuentes confiables y bases de datos oficiales (Infoleg, INDEC, BCRA).
    \item Score de probabilidad (0–1) acompañado de evidencia explicable.
\end{itemize}

\subsection{Limitaciones y exclusiones}

Quedan explícitamente fuera del alcance del prototipo:

\begin{itemize}
    \item Análisis de contenido multimedia (imágenes, videos, audios).
    \item Detección de deepfakes o manipulación de imágenes.
    \item Detección de campañas coordinadas o grafos de propagación.
    \item Integración con plataformas de mensajería (WhatsApp, Telegram).
    \item Soporte a otros idiomas además del español.
    \item Versiones para navegadores Firefox, Safari, Edge u otros.
    \item Aplicaciones móviles nativas (iOS/Android).
    \item Almacenamiento de datos personales de usuarios.
    \item Procesamiento en streaming a escala masiva.
\end{itemize}

\subsection{Estructura del documento}

Completar.
